% =========================================================================
% Capítulo 4: Cuerpos Rígidos - Rotación y Movimiento Compuesto
% =========================================================================
% Archivo: C4_00_introduccion.tex
% =========================================================================

\chapter{Cuerpo Rígido}

\section*{Introducción}

Un cuerpo rígido es un objeto sólido cuya distancia entre cualquier par de puntos permanece constante durante el movimiento. Puede pensarse como el límite continuo de un sistema de partículas (Capítulo 3) en el que las restricciones internas congelan las distancias relativas: de los infinitos grados de libertad sobreviven solo seis —tres de traslación y tres de rotación—. Este capítulo estudia cómo los cuerpos rígidos se comportan bajo fuerzas y torques, combinando traslación del centro de masa y rotación en torno a él.

\subsection*{Diferencias con Sistemas de Partículas}

\begin{itemize}
    \item Una distribución \textbf{continua} de masa en lugar de masas puntuales.
    \item La \textbf{rotación} es un aspecto dinámico integral del movimiento.
    \item El \textbf{momento de inercia} reemplaza a la masa en ecuaciones rotacionales.
\end{itemize}

\section{Propiedades Dinámicas de Cuerpos Rígidos}

\subsection*{Momento de Inercia}

El momento de inercia $I$ es una medida de la resistencia de un cuerpo a cambiar su rotación:

$$I = \int r^2 \, dm$$

donde $r$ es la distancia perpendicular del elemento de masa $dm$ al eje de rotación.

Para un eje que no pasa por el centro de masa:
$$I_{\text{eje}} = I_{\text{cm}} + Md^2 \quad \text{(Teorema de Ejes Paralelos)}$$

\subsection*{Tensor de Inercia}

Para rotaciones arbitrarias en 3D:

$$\mathbf{I} = \begin{pmatrix} I_{xx} & -I_{xy} & -I_{xz} \\ -I_{xy} & I_{yy} & -I_{yz} \\ -I_{xz} & -I_{yz} & I_{zz} \end{pmatrix}$$

\subsection*{Energía Cinética de Rotación}

Para un cuerpo rotando con velocidad angular $\vec{\omega}$:

$$T_{\text{rot}} = \frac{1}{2}\vec{\omega} \cdot \mathbf{I} \cdot \vec{\omega}$$

\subsection*{Momento Angular de un Cuerpo Rígido}

$$\vec{L} = \mathbf{I} \cdot \vec{\omega}$$

\section{Ecuaciones de Movimiento}

\subsection*{Movimiento de Traslación}

$$M\vec{a}_{\text{cm}} = \sum \vec{F}_{\text{ext}}$$

El centro de masa se mueve como si fuera una partícula de masa $M$ sometida a las fuerzas externas.

\subsection*{Movimiento de Rotación}

$$\frac{d\vec{L}}{dt} = \sum \vec{\tau}_{\text{ext}}$$

o, para rotación alrededor de un eje fijo:

$$I\alpha = \tau$$

donde $\alpha$ es la aceleración angular y $\tau$ es el torque total.

\subsection*{Rodadura sin Deslizar}

Para un cuerpo que rueda sin deslizar sobre una superficie:

$$v_{\text{cm}} = R\omega$$

$$a_{\text{cm}} = R\alpha$$

Esta restricción de movimiento relaciona traslación y rotación. Como muestra la Figura \ref{fig:c4_intro_rodadura}, el punto de contacto $C$ tiene velocidad instantánea nula, el centro $G$ se mueve con $v_{cm}=R\omega$ y el punto más alto del disco con $2R\omega$: el campo de velocidades del cuerpo rígido equivale a una rotación pura en torno a $C$.

% Figura: rodadura sin deslizar
\begin{figure}[!ht]
    \centering
    \begin{tikzpicture}[>=Latex, scale=1.5]

        % --- CONFIGURACIÓN DE PARÁMETROS ---
        \def\R{1.0}     % Radio del disco [u]
        \def\Ls{5.0}    % Largo del suelo [u]

        % --- SOPORTES Y ELEMENTOS FIJOS ---
        \draw[thick, gris1] (-\Ls/2,0) -- (\Ls/2,0);
        \fill[pattern=north east lines, pattern color=gris1!60] (-\Ls/2,-0.15) rectangle (\Ls/2,0);

        % --- GEOMETRÍA PRINCIPAL ---
        \draw[thick, fill=gris1!10] (0,\R) circle (\R);
        \coordinate (G) at (0,\R);
        \coordinate (C) at (0,0);
        \coordinate (T) at (0,2*\R);
        \draw[mydashed, gris1!60] (C) -- (T);

        % --- VECTORES Y FUERZAS ---
        % Campo de velocidades (rotación pura en torno a C)
        \draw[->, very thick, azul] (G) -- ++(1.0,0) node[right] {$v_{cm}=R\omega$};
        \draw[->, very thick, azul] (T) -- ++(2.0,0) node[right] {$2R\omega$};
        % Velocidad angular
        \draw[->, very thick, azul2] (0.55,\R+0.35) arc (60:-40:0.55) node[midway, right] {$\omega$};

        % --- ANOTACIONES Y ETIQUETAS ---
        \fill[gris1] (G) circle (1.3pt);
        \node[below left] at (G) {$G$};
        \fill[myred] (C) circle (1.5pt);
        \node[below right, myred] at (C) {$C$ \ ($\vec{v}_C=\vec{0}$)};

    \end{tikzpicture}
    \caption{Rodadura sin deslizar: campo de velocidades del disco}
    \label{fig:c4_intro_rodadura}
\end{figure}

\section{Ejemplos de Momentos de Inercia}

\begin{itemize}
    \item \textbf{Cilindro sólido}: $I = \frac{1}{2}MR^2$
    \item \textbf{Cilindro hueco}: $I = MR^2$
    \item \textbf{Esfera sólida}: $I = \frac{2}{5}MR^2$
    \item \textbf{Barra respecto al centro}: $I = \frac{1}{12}ML^2$
    \item \textbf{Barra respecto a un extremo}: $I = \frac{1}{3}ML^2$
\end{itemize}

\section*{Objetivos del Capítulo}

Al completar este capítulo, deberás:
\begin{itemize}
    \item Calcular momentos de inercia para geometrías complejas.
    \item Analizar movimiento combinado de traslación y rotación.
    \item Resolver problemas de rodadura y pivoteo.
    \item Aplicar conservación de momento angular.
    \item Entender sistemas de cuerpos rígidos acoplados.
\end{itemize}

\newpage
